\documentclass[12pt,a4paper]{article}
\usepackage{geometry}
\geometry{a4paper, margin=2cm}
\usepackage[utf8]{inputenc}
\usepackage[T1]{fontenc}
\usepackage[italian]{babel}
\usepackage{verbatim}
\usepackage{subcaption}
\usepackage{graphicx}
\usepackage{fancyhdr}
\usepackage{titlesec}
\usepackage{lmodern}
\usepackage{geometry}
\usepackage{setspace}
\usepackage{parskip}
\usepackage{listings}
\usepackage{xcolor}
\usepackage{enumitem}
\usepackage{float}
\usepackage{booktabs}
\usepackage{csquotes} 

\begin{document}

\begin{titlepage}
    \centering
    \vspace*{1.5cm}
    {\scshape\LARGE Università degli Studi di Udine \par}
    \vspace{0.4cm}
    {\scshape\Large Dipartimento di Scienze Matematiche, Informatiche e Fisiche\par}
    \vspace{0.4cm}
    {\large IBML - Internet of Things, Big Data, Machine Learning \par}
    \vspace{0.4cm}
    \vspace{1.5cm}
    
\includegraphics[width=0.25\textwidth]{immagini/Logo.pdf}\par
    \vspace{1cm}
    \rule{\textwidth}{0.5pt}\vspace{0.4cm}
    {\bfseries\LARGE Basi di dati\\[0.2cm]
    Progetto Data Base  - Gruppo 11 \par}
    \rule{\textwidth}{0.5pt}
    \vfill
    \begin{flushleft}
    \textit{\textbf{Autori}}\\
    Joshua Rumiato,  168062\hfill 168062@spes.uniud.it \\
    Kevin Polesello,  167947\hfill 167947@spes.uniud.it \\
    Mattia Nadal,  167372 \hfill 167372@spes.uniud.it \\
    Francesco Vignandel,   167024\hfill 167024@spes.uniud.it
    \end{flushleft}
\end{titlepage}

\newpage

\tableofcontents

\newpage

\section{Analisi dei requisiti}
\subsection{Traccia}
Il sistema deve gestire un campo-scuola in tutti gli aspetti organizzativi. Il campo è caratterizzato dai seguenti elementi:
\begin{itemize}
    \item \textbf{Edificio:} identificato dal relativo indirizzo.
    \item \textbf{Partecipanti:} divisi in \textit{Equipe} e \textit{Animati} (ragazzi).
\end{itemize}
Per ogni partecipante devono essere registrati i seguenti dati comuni:
\begin{itemize}
    \item Codice Fiscale e foto del documento.
    \item Nome, cognome, sesso (maschio/femmina).
    \item Data e luogo di nascita.
    \item Intolleranze alimentari.
    \item \textbf{Note aggiuntive:} patologie, allergie, medicinali (modalità e orari), situazioni particolari (es. bullismo) e variazioni logistiche (arrivi posticipati o partenze anticipate).
\end{itemize}
L'equipe è composta da cuochi, animatori, capi-campo e collaboratori. Per ogni membro si conservano i seguenti dati:
\begin{itemize}
    \item \textit{Recapiti} (numero di telefono ed email personale)
    \item \textbf{Valutazione finale:} un voto da 1 a 5 corredato da un commento testuale.
\end{itemize}
Per ogni ragazzo deve essere salvato il contatto (telefono ed email) di un genitore o tutore legale.
Tutti i partecipanti alloggiano nella struttura secondo le seguenti regole:
\begin{itemize}
    \item \textbf{Dati camera:} numero identificativo, numero di posti letto e piano dell'edificio.
    \item \textbf{Vincoli ragazzi:} le camere dei ragazzi non possono essere miste. Devono avere un animatore responsabile dello stesso sesso (es. animatrici per camere femminili).
    \item \textbf{Altro personale:} il resto dell'equipe e gli animatori non assegnati alle camere dei ragazzi possono dormire in camere comuni senza vincoli di sesso.
    \item \textbf{Capienza:} non è consentito superare il numero di posti letto previsti per ogni camera.
\end{itemize}
Gli animati sono divisi in squadre, ognuna definita da:
\begin{itemize}
    \item Nome, slogan e un animatore responsabile.
    \item Punteggio cumulativo per le attività di routine.
    \item Turni per le faccende domestiche.
\end{itemize}
Ogni attività ha un orario di inizio, un orario di fine e una descrizione. \textbf{Vincolo temporale:} deve esserci una sequenza rigida tra gli orari; le attività non possono sovrapporsi.
Le attività sono categorizzate come segue:
\begin{description}
    \item[PASTI:] Colazione, pranzo, merenda, cena. Sono consumati da tutti i partecipanti. I cuochi preparano i pasti con un responsabile specifico per ogni turno. Ogni pasto include una pietanza (nome, descrizione, allergeni).
    \item[FACCENDE:] Preparazione/spreparazione tavoli e pulizia stoviglie per i tre pasti principali.
    \item[ROUTINE:] Sveglia, attività fisica, attività di mattina/pomeriggio, docce, gioco libero, attività serale, verifica giornata. Ogni routine ha un nome, un responsabile e una valutazione testuale.
    \item[EXTRA:] Altre attività non categorizzate.
\end{description}
Il sistema tiene traccia dei materiali necessari per il campo. Ogni voce include:
\begin{itemize}
    \item Nome e quantità.
    \item Unità di misura e numero di utilizzi previsti.
    \item \textbf{Materiali specifici:} indicazione se il materiale è legato a una particolare attività di routine o se deve essere portato autonomamente dai partecipanti.
\end{itemize}

\subsection{Requisiti}
\subsection{Glossario}

\section{Progettazione concettuale}
\subsection{Modello Entità-Relazione}
\includegraphics[scale=0.5]{immagini/Diagramma-ER.pdf}

\section{Progettazione logica}
\subsection{Modello logico}
\subsection{Schema logico}

\section{Progettazione fisica}
\subsection{Definizione tabelle}
\subsection{Query SQL}

\section{Implementazione}
\subsection{Trigger}

\end{document}